\chapter{Validation}\label{ch:ml:validation}

A research team's model scores an \omterm{def:ml:why-financial-machine-learning-is-different:r2}{out-of-sample R-squared} of 1.3\% on the last seven years of a monthly stock panel,
three times what the best published models reach. The post-mortem finds nothing wrong with the model: the earnings
surprise it uses was joined to each stock by the end of the fiscal quarter, and the quarter's report reaches the market
a month later, in the very month whose return the model predicts. On the same data joined by filing date the model
scores $-0.19\%$, as it should: in the world of this chapter nothing is predictable at the moment of decision, so every
positive number below is a leak. Validation estimates how a model will do on data it has not seen; this chapter is
about the many ways the data it has not seen gets into it anyway, and how to catch each one.

\section{What validation estimates}

Validation answers two different questions, and a single number cannot answer both.

\begin{definition}[Model selection, hyperparameter search]\label{def:ml:validation:selection}
\emph{Model selection}\index{model selection} chooses, among candidate models or configurations, the one expected to
generalise best. A \emph{hyperparameter search}\index{hyperparameter search} is model selection over a set of
\omterm{def:ml:why-financial-machine-learning-is-different:sets}{hyperparameter} values: a grid, random draws (Bergstra and Bengio, 2012), or a sequential search that proposes new
values from the scores of the old.
\end{definition}

The second question is \emph{assessment}: how well will the chosen model do? The score that won the selection is a
biased answer to it.

\begin{proposition}[The winner's score is optimistic]\label{prop:ml:validation:max}
Let $\hat S_1,\dots,\hat S_K$ be unbiased estimates of the generalisation scores $S_1,\dots,S_K$ of $K$ candidates. Then
$\E[\max_k\hat S_k] \ge \max_kS_k$, with equality only if the maximum is attained by the same candidate on almost every
draw with no estimation error.
\end{proposition}

\begin{proof}
$\max$ is convex, so by Jensen's inequality $\E[\max_k\hat S_k] \ge \max_k\E[\hat S_k] = \max_kS_k$.
\end{proof}

The bias grows with the number of candidates and with the noise of each score, which on market data is large. Book~4
(chapter~12) turned this into the deflated Sharpe ratio; here it is the reason a selection score must never be reported
as an assessment.

\section{Splits that respect time}

The chapter's world is \texttt{firm.mlsynth}'s panel with a ceiling of zero: 200 stocks, 240 months, twenty
characteristics that carry nothing. Each stock reports quarterly; the report of the quarter ending at month $t$ is filed
during month $t+1$ and moves that month's return by 2\% per unit of standardised surprise. The pipeline is careful
everywhere else: gradient-boosted trees trained on returns demeaned by month, on the first 160 months less one, scored
against zero on the last 80.

Book~7 (chapter~20) built the splits that respect time: purged $k$-fold with an embargo, combinatorial purged
cross-validation and walk-forward analysis. \texttt{firm.cvsplit} wraps them as scikit-learn cross-validators, so the
library's search and scoring tools use them unchanged (\cref{fig:ml:validation:splits}). The label overlap they guard
against is severe on panels: with three-month labels sampled monthly and 4-fold cross-validation, a shuffled fold leaves
38\,080 training labels that share a month with a test label; the purged fold leaves none. Deep trees on the twenty
noise characteristics score 3.47\% with shuffled folds and $-3.37\%$ with purged ones. The characteristics are
persistent, so a flexible model can recognise a stock in a given year and read the neighbouring months' labels off it.
A classifier asked to tell the first half of the months from the second from the characteristics alone reaches an AUC
of 0.71: the features locate time, which is the warning.

\begin{omfigure}
\begin{tikzpicture}[x=0.72cm,y=0.55cm,font=\footnotesize]
\def\row#1#2{\node[anchor=east] at (-0.2,#1) {#2};}
\row{4}{shuffled $k$-fold}
\foreach \i/\c in {0/1,1/0,2/0,3/1,4/0,5/0,6/1,7/0,8/0,9/1,10/0,11/0,12/1,13/0,14/0,15/1} {
  \ifnum\c=1 \fill[omThm!45] (\i,3.7) rectangle ++(0.92,0.6); \else \fill[omDef!30] (\i,3.7) rectangle ++(0.92,0.6); \fi }
\row{3}{purged $k$-fold, embargo}
\foreach \i in {0,1,2,3,4,10,11,12,13,14,15} { \fill[omDef!30] (\i,2.7) rectangle ++(0.92,0.6); }
\foreach \i in {6,7,8} { \fill[omThm!45] (\i,2.7) rectangle ++(0.92,0.6); }
\fill[gray!25] (5,2.7) rectangle ++(0.92,0.6); \fill[gray!25] (9,2.7) rectangle ++(0.92,0.6);
\row{2}{walk-forward}
\foreach \i in {2,3,4,5,6,7,8,9} { \fill[omDef!30] (\i,1.7) rectangle ++(0.92,0.6); }
\foreach \i in {10,11} { \fill[omThm!45] (\i,1.7) rectangle ++(0.92,0.6); }
\row{1}{nested (one outer fold)}
\foreach \i in {0,1,2,3,4,5,6,7,8,9} { \fill[omDef!30] (\i,0.7) rectangle ++(0.92,0.6); }
\foreach \i in {3,4,5} { \fill[omProp!55] (\i,0.7) rectangle ++(0.92,0.6); }
\foreach \i in {12,13,14,15} { \fill[omThm!45] (\i,0.7) rectangle ++(0.92,0.6); }
\fill[gray!25] (10,0.7) rectangle ++(0.92,0.6); \fill[gray!25] (11,0.7) rectangle ++(0.92,0.6);
\draw[->] (0,0.1) -- (16.3,0.1) node[right]{time};
\fill[omDef!30] (1,-0.9) rectangle ++(0.5,0.4); \node[anchor=west] at (1.6,-0.7) {train};
\fill[omThm!45] (3.6,-0.9) rectangle ++(0.5,0.4); \node[anchor=west] at (4.2,-0.7) {test};
\fill[omProp!55] (6,-0.9) rectangle ++(0.5,0.4); \node[anchor=west] at (6.6,-0.7) {inner validation};
\fill[gray!25] (10.1,-0.9) rectangle ++(0.5,0.4); \node[anchor=west] at (10.7,-0.7) {purged or embargoed};
\end{tikzpicture}

\omcaption{Four ways to split sixteen periods. Shuffled folds scatter test periods among training periods, so
overlapping labels and persistent features leak; purging removes the training labels that overlap the test block and the
embargo those just after it; walk-forward trains only on the past; \omterm{def:ml:validation:nested}{nested cross-validation} chooses \omterm{def:ml:why-financial-machine-learning-is-different:sets}{hyperparameters} on
inner folds of the outer training data only.}\label{fig:ml:validation:splits}
\end{omfigure}

\section{Choosing hyperparameters without fooling yourself}

\begin{definition}[Nested cross-validation]\label{def:ml:validation:nested}
\emph{Nested cross-validation}\index{nested cross-validation} runs an inner cross-validation inside each \omterm{def:ml:why-financial-machine-learning-is-different:sets}{training set}
of an outer one: the inner folds choose the \omterm{def:ml:why-financial-machine-learning-is-different:sets}{hyperparameters}, the configuration so chosen is refitted on the outer
\omterm{def:ml:why-financial-machine-learning-is-different:sets}{training set} and scored on the outer test fold. The mean outer score assesses the whole procedure, selection included.
\end{definition}

Selection noise is easiest to see where there is nothing to select. Thirty runs of one configuration of boosted trees,
identical except for their random seed, are trained on the chapter's first 159 months. Their R-squared on months 160
to 199 ranges from $-0.60\%$ to $-0.27\%$, with a median of $-0.40\%$; the best of them looks 0.13 points better than
the median. On months 200 to 239, which the choice never saw, that run scores $-0.20\%$ and the median run $-0.19\%$
(\cref{fig:ml:validation:seeds}): across the thirty seeds the two periods' scores correlate at $-0.16$. The 0.13 points
were the selection's own noise, as \cref{prop:ml:validation:max} says.

\begin{omfigure}
\begin{tikzpicture}
\begin{axis}[width=8.6cm,height=6.2cm,xlabel={R-squared on the selection months (\%)},ylabel={R-squared on untouched months (\%)},
  tick label style={font=\small},label style={font=\small},xmin=-0.65,xmax=-0.2,ymin=-0.35,ymax=0.18,grid=major,
  grid style={gray!25},table/col sep=comma,
  xticklabel style={/pgf/number format/fixed,/pgf/number format/precision=1},scaled x ticks=false]
\addplot[only marks,mark=*,mark size=1.6pt,omDef] table[x=selection,y=fresh]{figdata/ml/03-validation/seeds.csv};
\draw[omThm,thick] (axis cs:-0.2723,-0.2) circle[radius=4pt];
\node[font=\footnotesize,omThm,anchor=south] at (axis cs:-0.2723,-0.165) {chosen};
\end{axis}
\end{tikzpicture}

\omcaption{Thirty seeds of the same boosted trees on a world with nothing to predict: the score on the months used to
choose among them against the score on months nobody looked at. The run chosen for its selection score (circled) is
ordinary afterwards. Data: \texttt{ml\_validation.selection}.}\label{fig:ml:validation:seeds}
\end{omfigure}

\omterm{def:ml:validation:nested}{Nested cross-validation} corrects the report. With eight such seeds as the candidates, three outer purged folds and three
inner ones on every other month (so labels do not overlap), the flat search reports its best mean outer score,
$-0.62\%$; the nested procedure, choosing inside each outer \omterm{def:ml:why-financial-machine-learning-is-different:sets}{training set}, scores $-0.85\%$. The honest number is the
worse one. The candidates were seeds, which makes the point plainly; with real \omterm{def:ml:why-financial-machine-learning-is-different:sets}{hyperparameters} the same bias sits on
top of whatever difference the \omterm{def:ml:why-financial-machine-learning-is-different:sets}{hyperparameters} make.

\begin{method}[A search that can be reported]\label{met:ml:validation:search}
\begin{enumerate}
\item Fix the splits (purged folds or walk-forward, from the labels' spans) and the score before the first run.
\item Search on inner folds only; record every configuration tried in the research log (Book~7, chapter~1): the trial
count is part of the result.
\item Report the nested (or walk-forward) score of the procedure, never the winning inner score.
\item Keep a final holdout that no search has touched, and use it once.
\end{enumerate}
\end{method}

\section{Leakage hunting}

\begin{definition}[Target leakage, train--test contamination]\label{def:ml:validation:leakage}
\emph{Target leakage}\index{target leakage} is information in a feature, or in any step of the pipeline, that would not
have been available at the decision time the prediction stands for, typically because it was derived from the target
or recorded after it. \emph{Train--test contamination}\index{train--test contamination} is information about the test
data reaching the fit through the data themselves: overlapping labels or duplicated records on both sides of a split,
or statistics computed on the whole sample.
\end{definition}

Book~7 named look-ahead bias (chapter~3) and survivorship bias; leakage is their general form, and the literature on
data mining calls it one of the commonest mistakes in the field (Kaufman, Rosset, Perlich and Stitelman). The chapter
plants five leaks in its careful pipeline (\cref{tab:ml:validation:leaks}). Two are in the features: the period-end join
of the surprise, and \emph{target encoding}, here the stock's industry's average return in the same month, a statistic
of the target that happens to be available in the table when the model is fitted. One is in the feature selection:
screening 1\,000 noise candidates on their information coefficient over the whole sample before the split, and keeping
ten. One is in the splits (shuffled folds on overlapping labels), and one in the choice (the best seed by test score).

\begin{table}[htbp]
\centering\small
\begin{tabular}{@{}lccl@{}}
\toprule
leak & leaky $R^2_{\mathrm{oos}}$ & honest version & caught by\\
\midrule
period-end join of the surprise & 1.27\% & $-0.19\%$ (filing date) & truncation test\\
target encoding, same month & 8.97\% & $-0.19\%$ (none) & canary, truncation test\\
screening on the whole sample & 0.13\% & $-0.16\%$ (training months) & canary, truncation test\\
shuffled folds, overlapping labels & 3.47\% & $-3.37\%$ (purged folds) & fold-overlap check\\
best of 30 seeds on the test & $-0.27\%$ & $-0.20\%$ (untouched months) & holdout, nested search\\
\bottomrule
\end{tabular}
\caption{Five leaks in a world where nothing is predictable at decision time (200 stocks, 240 months). Every positive
number is fake; the last row's leak is a gain of 0.13 points over the median seed that does not recur. Screening uses
ridge regression on the ten screened features, the others boosted trees. Data: \texttt{ml\_validation.leaks}.}\label{tab:ml:validation:leaks}
\end{table}

\begin{definition}[Adversarial validation, leakage audit]\label{def:ml:validation:audit}
\emph{Adversarial validation}\index{adversarial validation} trains a classifier to tell training rows from test rows; an
AUC well above one half means the features locate the split, so shuffled folds will leak and the test period differs
from the training period. A \emph{leakage audit}\index{leakage audit} is the set of checks a pipeline passes before its
score is believed: at least a truncation test of every feature, a canary run of the whole pipeline, a fold-overlap
check of the splits, and adversarial validation.
\end{definition}

The truncation test recomputes each feature on the data known at a date and compares it with the stored value: on the
surprise stored by filing month, the period-end join's feature differs by up to 4.77 standard deviations, the filing-date
join's by exactly zero; the screened features differ by 5.13, because the screen chose different candidates with less
data. The canary runs the whole pipeline on a target that is noise by construction, with the real target's shape: the
target-encoding pipeline still scores 4.03\% on noise (standard error 0.02\%) and the screening pipeline 0.17\% (0.03\%),
which no honest pipeline can. The canary does not catch the period-end join ($-0.44\%$), because that feature leaks
about the real target, not about the pipeline's own: the audit needs both tests. The fold-overlap check counts training
labels whose spans touch the test block; the holdout catches selection. No single check catches all five.

\begin{method}[The leakage audit]\label{met:ml:validation:audit}
\begin{enumerate}
\item Store every input with its knowledge time (Book~7's bitemporal store) and recompute each feature on the data
known at a sample of dates; any difference is a leak.
\item Run the whole pipeline on a canary target with the real target's structure (same overlap, same scale); a score
more than two standard errors above chance is a leak in the pipeline.
\item Check that no training label overlaps a test label; run \omterm{def:ml:validation:audit}{adversarial validation} on the features.
\item Keep a holdout outside every search, and count the trials.
\end{enumerate}
\end{method}

\section{Tutorial: five leaks}\label{tut:ml:validation:tut}

\begin{tutorial}
\textbf{Goal.} Plant five leaks in a pipeline on a world with nothing to predict, measure the R-squared each one fakes,
catch each with a detector, and compare flat and nested searches. \textbf{End state:} \cref{tab:ml:validation:leaks},
\cref{fig:ml:validation:seeds}, the detector numbers.
\begin{tutsteps}
\item \textbf{The period-end join}: one line decides whether the quarter ending at $t$ appears at $t$ or at $t+1$.
\omcode{code/ml/03-validation/python/ml_validation.py}{55}{63}{The surprise as the pipeline sees it, by filing date or by period end.}\label{lst:ml:validation:join}
\item \textbf{\omterm{def:ml:validation:nested}{Nested cross-validation}} over any scikit-learn estimator and any splitters.
\omcode{code/firm/cvsplit/firm_cvsplit.py}{70}{92}{Nested cross-validation and the flat search's optimistic best.}\label{lst:ml:validation:nested}
\item \textbf{The truncation test and the canary.}
\omcode{code/firm/cvsplit/firm_cvsplit.py}{102}{119}{Two detectors of the leakage audit.}\label{lst:ml:validation:detectors}
\item \textbf{Run} \texttt{leaks()}, \texttt{overlap\_detectors()}, \texttt{detectors()}, \texttt{nested()} and
\texttt{fig\_validation.py}.
\end{tutsteps}
\textbf{What to change next.} Replace the same-month target encoding by one computed on the previous month's returns and
check that the truncation test passes; raise the announcement effect to 4\% and see the period-end join's fake R-squared
grow with its square.
\end{tutorial}

\section{Build: splitters and leak detectors}\label{bld:ml:validation:cvsplit}

\begin{build}
\bfield{Purpose} Every model of the firm is validated with splits that respect time, selected without touching its
assessment, and passed through a \omterm{def:ml:validation:audit}{leakage audit}.
\bfield{Interface} \texttt{PurgedKFold(t0, t1, n\_splits, embargo)}, \texttt{CPCVSplit(t0, t1, n\_folds, n\_test,
embargo)}, \texttt{WalkForwardSplit(n, train, test, step, anchored)} (scikit-learn cross-validators on Book~7's
\texttt{firm.overfit}); \texttt{nested\_cv(make\_model, grid, X, y, outer, inner, score)};
\texttt{fold\_overlap}, \texttt{truncation\_test}, \texttt{canary}, \texttt{adversarial\_validation}.
\bfield{Rules} Rows in time order; every label carries its span; a score used to choose is never reported as an
assessment.
\bfield{Acceptance tests} \texttt{code/firm/cvsplit/tests/}: the splitters work inside \texttt{cross\_val\_score} and
purged folds leave no overlap; nested search picks the right ridge penalty and the flat best is never below the nested
score; the truncation test flags a centred moving average and passes a causal one; the canary flags a pipeline that
scores a feature made of its target; \omterm{def:ml:validation:audit}{adversarial validation} is near one half on one distribution and high on two.
\bfield{Stretch} A dependency tracer that records which rows each feature value read; a canary that permutes the target
in blocks as long as the labels.
\end{build}

\begin{omsources}
\item S.~Varma and R.~Simon, ``Bias in error estimation when using cross-validation for model selection'', \emph{BMC
Bioinformatics} 7, 2006.
\item G.~C.~Cawley and N.~L.~C.~Talbot, ``On over-fitting in model selection and subsequent selection bias in performance
evaluation'', \emph{Journal of Machine Learning Research} 11, 2010.
\item S.~Kaufman, S.~Rosset, C.~Perlich and O.~Stitelman, ``Leakage in data mining: formulation, detection, and
avoidance'', \emph{ACM Transactions on Knowledge Discovery from Data} 6(4), 2012.
\item J.~Bergstra and Y.~Bengio, ``Random search for hyper-parameter optimization'', \emph{Journal of Machine Learning
Research} 13, 2012.
\end{omsources}

\section{Exercises}

\begin{exercise}[$\star$]\label{exo:ml:validation:1}
Five \omterm{def:ml:why-financial-machine-learning-is-different:sets}{hyperparameter} settings have true \omterm{def:ml:why-financial-machine-learning-is-different:r2}{out-of-sample R-squared} of 0.30\% each, and each validation estimate has a
standard error of 0.10\%, independently. Roughly what does the best validation score average? (The expected maximum of
five standard normals is 1.16.)
\end{exercise}

\begin{exercise}[$\star$]\label{exo:ml:validation:2}
Labels span three months and are sampled monthly; the test fold is months 100 to 139. With purging and a one-month
embargo, which training months go?
\end{exercise}

\begin{exercise}[$\star$]\label{exo:ml:validation:3}
A canary run of a pipeline gives R-squared scores of 0.21\%, 0.15\% and 0.18\% on three noise targets. Is the pipeline
leaking? Compute the mean and its standard error.
\end{exercise}

\begin{exercise}[$\star\star$]\label{exo:ml:validation:4}
Why does the canary miss the period-end join, while the truncation test catches it? Which of the two catches target
encoding, and why both?
\end{exercise}

\begin{exercise}[$\star\star$]\label{exo:ml:validation:5}
\omterm{def:ml:validation:audit}{Adversarial validation} of a model's training and test months gives an AUC of 0.95. List two different conclusions this
supports and one it does not.
\end{exercise}

\begin{exercise}[$\star\star$]\label{exo:ml:validation:6}
\emph{Find the flaw.} ``We standardise every feature with the mean and standard deviation of the whole dataset, then run
purged cross-validation, so there is no leakage.''
\end{exercise}

\begin{exercise}[$\star\star\star$]\label{exo:ml:validation:7}
\emph{Coding.} In \texttt{ml\_validation}, set the announcement effect \texttt{ANN} to 4\% and rerun the period-end and
filing-date pipelines. By how much does the fake R-squared grow, and why about fourfold?
\end{exercise}

\begin{exercise}[$\star\star\star$]\label{exo:ml:validation:8}
Prove that \omterm{def:ml:validation:nested}{nested cross-validation}'s outer score is an unbiased estimate of the procedure's generalisation score when
the outer folds are independent of the procedure's choices, and explain why the flat best is not.
\end{exercise}

\section{Problem: Five Leaks}

\begin{problem}[{Weekend problem --- auditing a pipeline that cannot win}]\label{pb:ml:validation:1}
The chapter's world: 200 stocks, 240 months, characteristics that carry nothing, and earnings reports filed a month
after each quarter ends.

\medskip\noindent\textbf{Part I --- The honest pipeline.}
\begin{enumerate}
\item What does the careful pipeline score, and why is a negative number the right answer here?
\item Why is one month removed between the training and the test months?
\item What does the filing-date join put in month $t$'s row, and what does the period-end join put there?
\item What announcement effect makes the period-end join's fake R-squared 1.27\%?
\end{enumerate}

\medskip\noindent\textbf{Part II --- The leaks.}
\begin{enumerate}[resume]
\item Give the fake R-squared of the target encoding and explain its size.
\item What do screening on the whole sample and on the training months give?
\item How many training labels share a month with a shuffled test fold, and how many with a purged one?
\item What do deep trees score with shuffled and with purged folds, and what lets them exploit the overlap?
\end{enumerate}

\medskip\noindent\textbf{Part III --- Selection.}
\begin{enumerate}[resume]
\item What range do the thirty seeds' selection scores cover, and how far is the best from the median?
\item What does the chosen seed score on untouched months?
\item What do the flat and the nested searches report?
\item Why does \cref{prop:ml:validation:max} predict the gap?
\end{enumerate}

\medskip\noindent\textbf{Part IV --- The audit.}
\begin{enumerate}[resume]
\item Which detector catches each leak, with its number?
\item Why can no single detector catch all five?
\item What does an adversarial AUC of 0.71 on the characteristics mean for this pipeline?
\item State the \emph{named result}: the R-squared each leak fakes on a target that cannot be predicted, and the
detector that catches it.
\item Write the audit as a checklist a reviewer signs.
\item Which of the five leaks would survive a paper-trading period, and for how long?
\item How should the research log record a leak once found?
\item In one sentence: what is validation for?
\end{enumerate}
\end{problem}

\section{Interview questions}

\begin{interviewq}[$\star$ \iqroles{researcher, mle}]\label{iq:ml:validation:1}
What is data leakage? Give three examples from financial data.
\end{interviewq}

\begin{interviewq}[$\star\star$ \iqroles{mle}]\label{iq:ml:validation:2}
You tuned twenty \omterm{def:ml:why-financial-machine-learning-is-different:sets}{hyperparameter} settings by 5-fold cross-validation and report the best fold-average score. What is
wrong, and what do you report instead?
\end{interviewq}

\begin{interviewq}[$\star\star$ \iqroles{researcher}]\label{iq:ml:validation:3}
How would you test whether a feature pipeline has look-ahead, without reading its code?
\end{interviewq}

\begin{interviewq}[$\star\star$ \iqroles{researcher, mle}]\label{iq:ml:validation:4}
What is \omterm{def:ml:validation:audit}{adversarial validation} and what does it tell you about a train--test split?
\end{interviewq}

\begin{interviewq}[$\star\star$ \iqroles{researcher}]\label{iq:ml:validation:5}
Why is target encoding dangerous with time series, and how do you do it safely?
\end{interviewq}

\begin{interviewq}[$\star\star\star$ \iqroles{researcher}]\label{iq:ml:validation:6}
A model's cross-validated score is excellent and its paper-trading score is at zero. Walk through how you would find
out whether the cause is leakage, \omterm{def:ml:why-financial-machine-learning-is-different:generalisation}{overfitting} or a change in the market.
\end{interviewq}
